%! Vector-Valued Functions — Unit 4, Lesson 4.9 — Student Packet Cover Sheet
\documentclass[10pt]{article}
\usepackage{precalculus-article}
\usepackage{precalculus-boxes}
\usepackage{ltablex}
\keepXColumns

\begin{document}

% Full-bleed navy banner
\begin{tikzpicture}[remember picture, overlay]
  \fill[navy] (current page.north west)
    rectangle ([yshift=-0.9in]current page.north east);
\end{tikzpicture}
\vspace{-0.8in}

\begin{center}
  {\color{white}\textbf{\LARGE AP Precalculus}}\\[4pt]
  {\color{skymid}\large\textbf{Unit 4: Functions Involving Parameters, Vectors, and Matrices}}\\[6pt]
  {\color{white}\textbf{\large Lesson 4.9 \quad Vector-Valued Functions}}
\end{center}

\vspace{0.22in}
\namedateperiod
\vspace{0.18in}

\begin{learningtargetbox}
\small
\begin{enumerate}[leftmargin=1.5em, itemsep=5pt]
  \item I can represent the position of a particle as a vector-valued function
        $\vec{p}(t)=\langle x(t),\,y(t)\rangle$ and interpret $\|\vec{p}(t)\|$
        as the distance from the origin.
  \item I can represent velocity as a vector-valued function $\vec{v}(t)=\langle x'(t),\,y'(t)\rangle$,
        interpret the sign of each component as direction of motion, and compute
        speed as $\|\vec{v}(t)\|$.
\end{enumerate}
\end{learningtargetbox}

\vspace{0.14in}

\begin{tocbox}
\small
\renewcommand{\arraystretch}{1.5}
\begin{tabularx}{\linewidth}{@{} c l X r @{}}
\rowcolor{sky}
\textbf{\#} & \textbf{Component} & \textbf{Description} & \textbf{Score} \\
\midrule
1 & Warm-Up        & Spiral review: parametric evaluation, vector magnitude, average rate of change & \blank{1.2cm} \\
2 & Guided Notes   & Position vectors, velocity vectors, speed, and worked examples & \blank{1.2cm} \\
3 & Group Activity & Tiered practice with position and velocity vector-valued functions & \blank{1.2cm} \\
4 & Exit Ticket    & Express position as a vector, find magnitude, compute speed & \blank{1.2cm} \\
5 & Homework       & Independent practice including an AP-style multiple-choice problem & \blank{1.2cm} \\
\midrule
  & \multicolumn{2}{r}{\textbf{Total}} & \blank{1.2cm} \\
\end{tabularx}
\end{tocbox}

\end{document}
